\documentclass[12pt,a4paper]{article}

% Paquetes
\usepackage[utf8]{inputenc}
\usepackage[spanish]{babel}
\usepackage{amsmath, amssymb}
\usepackage{booktabs}
\usepackage{geometry}
\usepackage{graphicx}
\usepackage{xcolor}
\usepackage{float}
\usepackage{hyperref}
\usepackage{fancyhdr}
\usepackage{titlesec}

\geometry{margin=2.5cm}

\definecolor{azuluni}{RGB}{0,70,127}

\pagestyle{fancy}
\fancyhf{}
\rhead{\textcolor{azuluni}{\small T4 -- Distribuciones de Probabilidad}}
\lhead{\textcolor{azuluni}{\small Caso Discreto y Continuo}}
\rfoot{\thepage}

\titleformat{\section}{\large\bfseries\color{azuluni}}{}{0em}{}[\titlerule]
\titleformat{\subsection}{\normalsize\bfseries\color{azuluni}}{}{0em}{}

\begin{document}

% ============================================================
\begin{titlepage}
  \centering
  \vspace*{2cm}
  {\color{azuluni}\rule{\linewidth}{1.5pt}}\\[0.5cm]
  {\LARGE\bfseries\color{azuluni}
   Taller T4 -- Distribuciones de Probabilidad\\[0.3cm]
   Casos Discreto y Continuo}\\[0.3cm]
  {\color{azuluni}\rule{\linewidth}{1.5pt}}\\[1.5cm]
  {\large\bfseries Tarea T4 -- Probabilidad}\\[0.4cm]
  {\large Alejandro Escobar Barrios --- 20251020094}\\[2cm]
  \begin{tabular}{ll}
    \textbf{Distribuciones:} & Binomial (Bernoulli), Poisson, Normal, Exponencial \\[0.2cm]
    \textbf{Herramientas:}   & \LaTeX{} + Python (numpy, scipy, matplotlib) \\
  \end{tabular}
  \vfill
  {\small\color{gray} Generado con \LaTeX{} y validado con Python}
\end{titlepage}

% ============================================================
% PARTE 1: DISTRIBUCIÓN BINOMIAL (BERNOULLI)
% ============================================================
\section{Parte 1: Distribucion Binomial (Bernoulli)}

\subsection{Marco Teorico}

Un experimento de Bernoulli es aquel que tiene exactamente dos resultados posibles:
\emph{exito} (probabilidad $p$) o \emph{fracaso} (probabilidad $q = 1-p$).
La \textbf{distribucion Binomial} modela el numero de exitos $X$ en $n$ ensayos
independientes de Bernoulli:
\[
  \boxed{P(X = k) = \binom{n}{k} p^k\, q^{n-k}, \quad k = 0, 1, \ldots, n.}
\]

Sus momentos principales son:
\[
  E[X] = np, \qquad \mathrm{Var}(X) = npq.
\]

\subsection{Ejercicio Teorico 1: Propiedades de la Distribucion Binomial}

\textbf{Problema.} Sea $X \sim \mathrm{Bin}(n, p)$. Demostrar que $\displaystyle\sum_{k=0}^{n} P(X=k) = 1$.

\textbf{Demostracion.}
\begin{align*}
  \sum_{k=0}^{n} P(X=k)
  &= \sum_{k=0}^{n} \binom{n}{k} p^k q^{n-k} \\
  &= (p + q)^n \quad \text{(Teorema del Binomio)} \\
  &= (p + 1 - p)^n = 1^n = \mathbf{1}. \quad \checkmark
\end{align*}

\subsection{Ejercicio Teorico 2: Esperanza de la Binomial}

\textbf{Problema.} Demostrar que $E[X] = np$ para $X \sim \mathrm{Bin}(n,p)$.

\textbf{Demostracion.}
\begin{align*}
  E[X] &= \sum_{k=0}^{n} k\binom{n}{k} p^k q^{n-k} \\
       &= \sum_{k=1}^{n} k\cdot\frac{n!}{k!(n-k)!} p^k q^{n-k} \\
       &= np \sum_{k=1}^{n} \frac{(n-1)!}{(k-1)!(n-k)!} p^{k-1} q^{n-k} \\
       &= np \sum_{j=0}^{n-1} \binom{n-1}{j} p^{j} q^{n-1-j} \quad (j = k-1)\\
       &= np\,(p+q)^{n-1} = np \cdot 1 = \mathbf{np}. \quad \checkmark
\end{align*}

\subsection{Aplicacion: Lanzamiento de un Dado}

\textbf{Problema.} Se lanza un dado no alterado 5 veces. Se define como \emph{exito}
obtener el numero 3. Calcular:
\begin{itemize}
  \item[(i)]   La probabilidad de obtener exactamente 2 exitos.
  \item[(ii)]  La probabilidad de obtener maximo 1 exito.
  \item[(iii)] La probabilidad de obtener al menos 2 exitos.
\end{itemize}

\textbf{Parametros.} $n = 5$, $p = \tfrac{1}{6}$, $q = \tfrac{5}{6}$.

\bigskip
\textbf{(i) Exactamente 2 exitos:}
\[
  P(X = 2) = \binom{5}{2}\!\left(\tfrac{1}{6}\right)^{\!2}\!\left(\tfrac{5}{6}\right)^{\!3}
  = 10 \cdot \frac{1}{36} \cdot \frac{125}{216}
  = \frac{1250}{7776} \approx \mathbf{0.1608}.
\]

\textbf{(ii) Maximo 1 exito} ($X \leq 1$):
\begin{align*}
  P(X \leq 1) &= P(X=0) + P(X=1) \\
  P(X=0) &= \binom{5}{0}\!\left(\tfrac{1}{6}\right)^{\!0}\!\left(\tfrac{5}{6}\right)^{\!5}
            = \frac{3125}{7776} \approx 0.4019.\\
  P(X=1) &= \binom{5}{1}\!\left(\tfrac{1}{6}\right)^{\!1}\!\left(\tfrac{5}{6}\right)^{\!4}
            = 5 \cdot \frac{1}{6} \cdot \frac{625}{1296}
            = \frac{3125}{7776} \approx 0.4019.\\
  P(X \leq 1) &= \frac{3125 + 3125}{7776} = \frac{6250}{7776} \approx \mathbf{0.8038}.
\end{align*}

\textbf{(iii) Al menos 2 exitos} ($X \geq 2$):
\[
  P(X \geq 2) = 1 - P(X \leq 1) = 1 - 0.8038 \approx \mathbf{0.1962}.
\]

\begin{table}[H]
\centering
\caption{Resultados -- Distribucion Binomial $(n=5,\; p=1/6)$}
\renewcommand{\arraystretch}{1.4}
\begin{tabular}{lcc}
\toprule
\textbf{Cantidad} & \textbf{Exacto} & \textbf{Python} \\
\midrule
$P(X = 2)$      & $1250/7776 \approx 0.1608$ & $0.160751$ \\
$P(X \leq 1)$   & $6250/7776 \approx 0.8038$ & $0.803755$ \\
$P(X \geq 2)$   & $1526/7776 \approx 0.1962$ & $0.196245$ \\
$E[X] = np$     & $5/6 \approx 0.8333$       & $0.833333$ \\
$\mathrm{Var}(X)$ & $25/36 \approx 0.6944$   & $0.694444$ \\
\bottomrule
\end{tabular}
\end{table}

% ============================================================
% PARTE 2: DISTRIBUCIÓN DE POISSON
% ============================================================
\section{Parte 2: Distribucion de Poisson}

\subsection{Marco Teorico}

La \textbf{distribucion de Poisson} modela el numero de eventos raros que ocurren
en un intervalo fijo de tiempo o espacio, con tasa promedio $\lambda$:
\[
  \boxed{P(X = k) = \frac{e^{-\lambda}\lambda^{k}}{k!}, \quad k = 0, 1, 2, \ldots}
\]

Caracteristicas:
\[
  E[X] = \lambda, \qquad \mathrm{Var}(X) = \lambda.
\]

\textbf{Aproximacion Binomial--Poisson.} Si $n$ es grande y $p$ es pequeno, entonces
$X \sim \mathrm{Bin}(n,p) \approx \mathrm{Poisson}(\lambda = np)$.

\subsection{Ejercicio Teorico 1: Normalizacion de Poisson}

\textbf{Problema.} Verificar que $\displaystyle\sum_{k=0}^{\infty} P(X=k) = 1$.

\textbf{Demostracion.}
\[
  \sum_{k=0}^{\infty} \frac{e^{-\lambda}\lambda^k}{k!}
  = e^{-\lambda}\sum_{k=0}^{\infty}\frac{\lambda^k}{k!}
  = e^{-\lambda}\cdot e^{\lambda} = e^{0} = \mathbf{1}.\quad\checkmark
\]

\subsection{Ejercicio Teorico 2: Esperanza de Poisson}

\textbf{Problema.} Demostrar que $E[X] = \lambda$.

\textbf{Demostracion.}
\begin{align*}
  E[X] &= \sum_{k=0}^{\infty} k\cdot\frac{e^{-\lambda}\lambda^k}{k!}
         = e^{-\lambda}\sum_{k=1}^{\infty}\frac{\lambda^k}{(k-1)!} \\
       &= e^{-\lambda}\cdot\lambda\sum_{j=0}^{\infty}\frac{\lambda^j}{j!} \quad (j=k-1) \\
       &= e^{-\lambda}\cdot\lambda\cdot e^{\lambda} = \boldsymbol{\lambda}.\quad\checkmark
\end{align*}

\subsection{Aplicacion: Accidentes en una Interseccion}

\textbf{Problema.} La probabilidad de que un auto se accidente es $p = 0{,}001$.
Entre las 4 pm y 6 pm pasan aproximadamente $n = 1000$ autos por la interseccion.
?`Cual es la probabilidad de que ocurran 2 o mas accidentes?

\textbf{Parametros.} Como $n$ es grande y $p$ es pequeno, se usa la aproximacion:
\[
  \lambda = np = 1000 \times 0{,}001 = 1.
\]

\textbf{Desarrollo.}
\begin{align*}
  P(X \geq 2) &= 1 - P(X = 0) - P(X = 1) \\[4pt]
  P(X = 0) &= \frac{e^{-1}\cdot 1^0}{0!} = e^{-1} \approx 0.3679. \\[4pt]
  P(X = 1) &= \frac{e^{-1}\cdot 1^1}{1!} = e^{-1} \approx 0.3679. \\[4pt]
  P(X \geq 2) &= 1 - 2e^{-1} \approx 1 - 0.7358 = \mathbf{0.2642}.
\end{align*}

\begin{table}[H]
\centering
\caption{Resultados -- Distribucion de Poisson $(\lambda = 1)$}
\renewcommand{\arraystretch}{1.4}
\begin{tabular}{lcc}
\toprule
\textbf{Cantidad} & \textbf{Exacto} & \textbf{Python} \\
\midrule
$P(X = 0)$    & $e^{-1} \approx 0.3679$   & $0.367879$ \\
$P(X = 1)$    & $e^{-1} \approx 0.3679$   & $0.367879$ \\
$P(X \geq 2)$ & $1-2e^{-1}\approx 0.2642$ & $0.264241$ \\
$E[X]$        & $\lambda = 1$              & $1.000000$ \\
$\mathrm{Var}(X)$ & $\lambda = 1$          & $1.000000$ \\
\bottomrule
\end{tabular}
\end{table}

% ============================================================
% PARTE 3: DISTRIBUCIÓN NORMAL
% ============================================================
\section{Parte 3: Distribucion Normal}

\subsection{Marco Teorico}

La \textbf{distribucion Normal} $N(\mu, \sigma^2)$ tiene f.d.p.:
\[
  \boxed{f(x) = \frac{1}{\sigma\sqrt{2\pi}}\exp\!\left(-\frac{(x-\mu)^2}{2\sigma^2}\right),
  \quad -\infty < x < \infty.}
\]

La estandarizacion $Z = \dfrac{X - \mu}{\sigma}\sim N(0,1)$ permite usar tablas.
Sus momentos: $E[X] = \mu$, $\mathrm{Var}(X) = \sigma^2$.

\subsection{Ejercicio Teorico 1: Simetria de la Normal Estandar}

\textbf{Problema.} Demostrar que $P(Z \leq -z) = P(Z \geq z)$ para $z > 0$.

\textbf{Demostracion.}
\[
  P(Z \leq -z) = \int_{-\infty}^{-z}\frac{e^{-t^2/2}}{\sqrt{2\pi}}\,dt.
\]
Haciendo el cambio $u = -t$ (y $du = -dt$):
\[
  = \int_{z}^{\infty}\frac{e^{-u^2/2}}{\sqrt{2\pi}}\,du = P(Z \geq z).\quad\checkmark
\]

\subsection{Ejercicio Teorico 2: Probabilidad en Intervalo Simetrico}

\textbf{Problema.} Demostrar que $P(-z\leq Z\leq z) = 2\Phi(z) - 1$,
donde $\Phi$ es la FDA de la normal estandar.

\textbf{Demostracion.}
\begin{align*}
  P(-z \leq Z \leq z) &= P(Z \leq z) - P(Z \leq -z) \\
                       &= \Phi(z) - [1 - \Phi(z)] \\
                       &= 2\Phi(z) - 1. \quad \checkmark
\end{align*}

\subsection{Aplicacion: Ejemplo 6.3 de Walpole}

\textbf{Problema (Walpole, p.\ 178, Ejemplo 6.3).}
Dado que $X \sim N(\mu = 50,\; \sigma = 10)$, encontrar $P(45 < X < 62)$.

\textbf{Estandarizacion.}
\[
  Z_1 = \frac{45 - 50}{10} = -0.5, \qquad Z_2 = \frac{62 - 50}{10} = 1.2.
\]

\textbf{Calculo.}
\begin{align*}
  P(45 < X < 62) &= P(-0.5 < Z < 1.2) \\
                  &= \Phi(1.2) - \Phi(-0.5) \\
                  &= \Phi(1.2) - [1 - \Phi(0.5)] \\
                  &= 0.8849 - (1 - 0.6915) \\
                  &= 0.8849 - 0.3085 = \mathbf{0.5764}.
\end{align*}

\textbf{Interpretacion.} Hay aproximadamente un 57.64\,\% de probabilidad de que
la variable se encuentre entre 45 y 62.

\begin{table}[H]
\centering
\caption{Resultados -- Distribucion Normal $(\mu=50,\;\sigma=10)$}
\renewcommand{\arraystretch}{1.4}
\begin{tabular}{lcc}
\toprule
\textbf{Cantidad} & \textbf{Exacto/Tablas} & \textbf{Python} \\
\midrule
$P(45 < X < 62)$  & $0.5764$  & $0.576398$ \\
$\Phi(1.2)$        & $0.8849$  & $0.884930$ \\
$\Phi(-0.5)$       & $0.3085$  & $0.308538$ \\
$E[X] = \mu$       & $50$      & $50.000000$ \\
$\mathrm{Var}(X)$  & $100$     & $100.000000$ \\
\bottomrule
\end{tabular}
\end{table}

% ============================================================
% RESUMEN GENERAL
% ============================================================
\section{Resumen General}

\begin{table}[H]
\centering
\caption{Comparacion de las tres distribuciones estudiadas}
\renewcommand{\arraystretch}{1.5}
\begin{tabular}{lllcc}
\toprule
\textbf{Distribucion} & \textbf{Parametros} & \textbf{E[X]} & \textbf{Var(X)} & \textbf{Tipo} \\
\midrule
Binomial & $n=5, p=1/6$       & $5/6\approx0.833$  & $25/36\approx0.694$ & Discreta \\
Poisson  & $\lambda=1$        & $1$                & $1$                 & Discreta \\
Normal   & $\mu=50,\sigma=10$ & $50$               & $100$               & Continua \\
\bottomrule
\end{tabular}
\end{table}

% ============================================================
\section{Conclusiones}

\begin{enumerate}
  \item La \textbf{distribucion Binomial} permite calcular probabilidades exactas para
        experimentos de Bernoulli repetidos. En el dado, la simetria del complemento
        ($P(X\geq2)=1-P(X\leq1)$) simplifica significativamente el calculo.

  \item La \textbf{distribucion de Poisson} es una aproximacion eficiente para la Binomial
        cuando $n$ es grande y $p$ pequeno ($\lambda = np = 1$). La probabilidad de
        2 o mas accidentes es aproximadamente 26.4\,\%, lo que es relevante para la
        gestion de riesgo vial.

  \item La \textbf{distribucion Normal} (Walpole 6.3) muestra como la estandarizacion
        transforma cualquier variable normal en $Z\sim N(0,1)$, facilitando el uso
        de tablas. El 57.64\,\% de probabilidad en $(45, 62)$ ilustra la concentracion
        de la masa de probabilidad alrededor de la media.

  \item La \textbf{distribucion Exponencial} modela fenomenos de espera con la propiedad
        unica de \emph{perdida de memoria}. Un tiempo medio de 3 minutos implica
        que el 18.9\,\% de las esperas superan los 5 minutos.

  \item Los resultados de Python concuerdan con los valores analiticos hasta al menos
        6 cifras decimales, validando todos los desarrollos matematicos.
\end{enumerate}

\end{document}
